\chapter{Fourierreeksen}\label{ch:b2:fourier}

Kan elk periodiek signaal uit zuivere sinussen en cosinussen worden
heropgebouwd? Fouriers vermetele ``ja'' bracht een eeuw analyse
voort. Dit hoofdstuk bewijst de twee pijlers die op dit niveau
bereikbaar zijn: de \emph{stelling van Dirichlet} (\omterm{def:b2:funcseq:def}{puntsgewijze}
reconstructie voor stuksgewijs $C^1$-functies, via de
\omterm{lem:b2:fourier:kernel}{Dirichlet-kern}) en de \emph{identiteit van Parseval} (de energie
van een signaal is de som van de energieën van zijn harmonischen),
en oogst de klassieke numerieke reeksen --- de som van Basel $\sum
1/n^2 = \pi^2/6$ voorop.

Overal zijn de functies $2\pi$-periodiek, stuksgewijs \omterm{def:b2:metric:continuity}{continu} en
complexwaardig; $\mathcal{C}$ duidt de \omterm{def:b2:metric:continuity}{continue} functies aan.

\section{Fouriercoëfficiënten}

\begin{definition}\label{def:b2:fourier:coefficients}
De \emph{Fouriercoëfficiënten}\index{Fouriercoëfficiënten} van $f$
zijn
\[
c_n(f) = \frac{1}{2\pi}\int_{-\pi}^{\pi} f(t)\,\eu^{-\iu n t}\,\dd
t \qquad (n \in \Z),
\]
en de coëfficiënten in reële vorm zijn $a_n = c_n + c_{-n}$, $b_n =
\iu(c_n - c_{-n})$, zodat de
\emph{Fourierpartiaalsommen}\index{Fourierpartiaalsommen} gegeven
worden door
\[
S_N(f)(t) = \sum_{n=-N}^{N} c_n(f)\,\eu^{\iu nt}
= \frac{a_0}{2} + \sum_{n=1}^{N}\bigl(a_n\cos nt + b_n \sin
nt\bigr).
\]
Definieer op $\mathcal{C}$ het \omterm{def:b2:hermitian:adjoint}{hermitische} inproduct $\langle f,
g\rangle = \frac{1}{2\pi}\int_{-\pi}^{\pi}\conj f\,g$: de
exponentiëlen $e_n(t) = \eu^{\iu nt}$ zijn \emph{orthonormaal}
($\langle e_m, e_n\rangle = \delta_{mn}$, rechtstreekse
berekening), en $c_n(f) = \langle e_n, f\rangle$: Fourieranalyse is
\omterm{def:b2:hermitian:adjoint}{hermitische} meetkunde (\cref{ch:b2:hermitian}) in oneindige
dimensie.
\end{definition}

\begin{proposition}[Ongelijkheid van Bessel]\label{prop:b2:fourier:bessel}
$S_N(f)$ is de orthogonale projectie van $f$ op de ruimte
$\mathcal{T}_N$ van de trigonometrische veeltermen van graad $\leq
N$, en\index{ongelijkheid van Bessel}
\[
\sum_{n=-N}^{N} \abs{c_n(f)}^2 \leq \norm f_2^2 =
\frac{1}{2\pi}\int_{-\pi}^{\pi} \abs f^2 :
\]
de reeks $\sum \abs{c_n}^2$ convergeert, en $c_n(f) \to 0$ als
$\abs n \to \infty$ (Riemann--Lebesgue voor coëfficiënten).
\end{proposition}

\begin{proof}
$f - S_N(f)$ staat loodrecht op elke $e_k$ met $\abs k \leq N$
($\langle e_k, f - S_N f\rangle = c_k - c_k = 0$): $S_Nf$ is de
orthogonale projectie op $\mathcal{T}_N =
\operatorname{Vect}(e_{-N}, \dots, e_N)$ (de projectiestelling uit
het volume van bachelorjaar 1, woordelijk in het \omterm{def:b2:hermitian:adjoint}{hermitische}
kader). Pythagoras: $\norm f_2^2 = \norm{S_Nf}_2^2 + \norm{f -
S_Nf}_2^2 \geq \norm{S_Nf}_2^2 = \sum_{\abs n \leq N}
\abs{c_n}^2$; laat $N \to \infty$.
\end{proof}

\begin{example}[Beste benadering, gemeten]\label{ex:b2:fourier:bestapprox}
Hoe goed benaderen trigonometrische veeltermen van lage graad de
zaagtand $f(t) = t$ (op $\intoo{-\pi}{\pi}$) in kwadratisch
gemiddelde? Volgens \cref{prop:b2:fourier:bessel} \emph{is} de
beste benadering van graad $N$ juist $S_N(f)$, met kwadratische
fout
\[
\norm{f - S_Nf}_2^2 = \norm f_2^2 -
\sum_{\abs n\leq N}\abs{c_n}^2 .
\]
Hier is $\norm f_2^2 = \frac{1}{2\pi}\int_{-\pi}^\pi t^2\dd t =
\frac{\pi^2}{3}$, en uit $b_n = \frac{2(-1)^{n+1}}{n}$
(\cref{ex:b2:fourier:basel}) volgt $\abs{c_n}^2 + \abs{c_{-n}}^2 =
\frac{b_n^2}{2} = \frac{2}{n^2}$. Dus
\[
\norm{f - S_Nf}_2^2
= \frac{\pi^2}{3} - \sum_{n=1}^{N}\frac{2}{n^2}
\qquad\text{: numeriek } 1.29,\ 0.79,\ 0.57,\ 0.44
\]
voor $N = 1, 2, 3, 4$ --- dalend, maar traag: de staart
$\sum_{n>N}\frac2{n^2} \sim \frac2N$ wordt bestuurd door het trage
verval $\frac1n$ van de coëfficiënten, dat zelf het handschrift
van de sprong is (\cref{exo:b2:fourier:6} achterstevoren gelezen).
Het inzicht om te onthouden: Parseval maakt van de kwaliteit van
een benadering de staart van een numerieke reeks --- en voorspelt,
nog vóór enige tekening, dat sprongen \omterm{def:b2:fourier:coefficients}{Fourierreeksen} met tegenzin
doen convergeren.
\end{example}

\begin{method}[Fouriercoëfficiënten efficiënt berekenen]\label{met:b2:fourier:coefficients}
Vóór je ook maar iets integreert:
\begin{enumerate}
  \item \emph{Pariteit:} een even $f$ heeft $b_n = 0$, een oneven
        $f$ heeft $a_n = 0$, en de overblijvende integralen
        herleiden tot $\frac2\pi \int_0^\pi$ --- half zoveel werk,
        dubbel zo betrouwbaar.
  \item \emph{Trigonometrische veeltermen zijn al klaar:}
        lineariseer producten ($\cos^3$, $\sin^2\cos$, \dots) en
        lees de coëfficiënten af (\cref{exo:b2:fourier:9}); de
        orthonormaliteit maakt elke verdere integratie
        overbodig.
  \item \emph{Complexe exponentiëlen voor exponentiëlen:} voor
        factoren $\eu^{at}$ of gedempte trillingen bereken je
        $c_n$ rechtstreeks --- één integraal van $\eu^{(a - \iu
        n)t}$ verslaat twee partiële integraties
        (\cref{exo:b2:fourier:10}).
  \item \emph{Differentieer een bekende ontwikkeling:} zijn de
        coëfficiënten van $f'$ bekend en is $f$ \omterm{def:b2:metric:continuity}{continu}, dan geeft
        $c_n(f) = \frac{c_n(f')}{\iu n}$ ($n \neq 0$) alles behalve
        $c_0$, dat het gemiddelde is --- vaak de snelste weg, en
        precies onder de hypothesen van
        \cref{thm:b2:fourier:parseval}~(1) geoorloofd.
\end{enumerate}
\end{method}

\section{De stelling van Dirichlet}

\begin{lemma}[Dirichlet-kern]\label{lem:b2:fourier:kernel}
$S_N(f)(x) = \frac{1}{2\pi}\int_{-\pi}^{\pi} f(x + u)\,D_N(u)\,\dd
u$, waarbij\index{Dirichlet-kern}
\[
D_N(u) = \sum_{n=-N}^{N} \eu^{\iu nu}
= \frac{\sin\bigl((N + \frac12)u\bigr)}{\sin\frac u2}
\quad (u \notin 2\pi\Z),
\qquad
\frac{1}{2\pi}\int_{-\pi}^{\pi} D_N = 1 .
\]
\end{lemma}

\begin{proof}
Vul de definitie van $c_n$ in $S_N$ in en verwissel som en
integraal (geoorloofd: de som is eindig):
\[
S_N(f)(x)
= \sum_{n=-N}^{N}\Bigl(\frac{1}{2\pi}\int_{-\pi}^{\pi}
f(t)\,\eu^{-\iu nt}\dd t\Bigr)\eu^{\iu nx}
= \frac{1}{2\pi}\int_{-\pi}^{\pi} f(t)
\sum_{n=-N}^{N}\eu^{\iu n(x-t)}\,\dd t ;
\]
substitueer $u = t - x$ en schuif het integratiesegment over
$2\pi$ terug naar $\intcc{-\pi}{\pi}$ dankzij de
$2\pi$-periodiciteit van de integrand; het \omterm{def:b2:quadratic:adjoint}{symmetrische}
indexbereik maakt $\sum_n\eu^{-\iu nu} = D_N(u)$. De gesloten
vorm: meetkundige som met reden $\eu^{\iu u}$,
\[
D_N(u) = \eu^{-\iu Nu}\,\frac{\eu^{\iu(2N+1)u} - 1}{\eu^{\iu u} -
1}
= \frac{\eu^{\iu(N + \frac12)u} - \eu^{-\iu(N+\frac12)u}}
{\eu^{\iu u/2} - \eu^{-\iu u/2}} ,
\]
wat het sinusquotiënt is. Haar gemiddelde is $1$: alleen $n = 0$
draagt bij.
\end{proof}

\begin{theorem}[Lemma van Riemann--Lebesgue]\label{thm:b2:fourier:riemannlebesgue}
Voor stuksgewijs \omterm{def:b2:metric:continuity}{continue} $g$ op een segment geldt $\int_a^b
g(t)\sin(\lambda t + \varphi)\,\dd t \to 0$ als $\lambda \to
+\infty$.\index{lemma van Riemann--Lebesgue}
\end{theorem}

\begin{proof}
Benader $g$ uniform door trapfuncties
(\cref{thm:b2:funcseq:weierstrass} is niet nodig --- de
elementaire benadering van stuksgewijs \omterm{def:b2:metric:continuity}{continue} functies door
trapfuncties volstaat) en integreer elke trap expliciet: elk stuk
draagt $O\bigl(\frac1\lambda\bigr)$ bij, en de benaderingsfout
draagt $\varepsilon(b - a)$ bij. Dit argument werd \omterm{def:b2:metric:complete}{volledig}
uitgevoerd als laatste oefening van het integratiehoofdstuk in het
volume van bachelorjaar 1; voor $C^1$-stukken kan men in plaats
daarvan partieel integreren en door $\frac C\lambda$ begrenzen.
\end{proof}

\begin{example}[Hoe snel sterven de coëfficiënten?]\label{ex:b2:fourier:decaytable}
Riemann--Lebesgue zegt dat de coëfficiënten naar $0$ streven; hun
\emph{snelheid} is een gladheidsmeter. Drie exemplaren uit dit
hoofdstuk en zijn oefeningen:
\[
\text{blokgolf: } b_n = \frac{4}{\pi n}\ (n\ \text{oneven}),
\qquad
\abs t : a_n = \frac{-4}{\pi n^2}\ (n\ \text{oneven}),
\qquad
\abs{\sin t} : a_{2k} = \frac{-4}{\pi(4k^2-1)} .
\]
Een sprong in $f$ (blokgolf, zaagtand) laat coëfficiënten van de
orde $\frac1n$ na: geen \omterm{def:b2:funcseq:series}{normale convergentie}, en een
Gibbs-overshoot bij de sprongen. \omterm{def:b2:metric:continuity}{Continuïteit} met een knik --- een
sprong in alleen $f'$ --- verbetert de orde tot $\frac{1}{n^2}$:
\omterm{def:b2:funcseq:series}{normale convergentie}, \omterm{def:b2:funcseq:def}{uniforme} reconstructie. In het algemeen
kopen $k$ afgeleiden $c_n = O(n^{-k})$ (\cref{exo:b2:fourier:6}),
en omgekeerd dwingt een spectrum dat sneller dan elke macht
vervalt af dat $f$ van klasse $C^\infty$ is (differentieer nu wél
geoorloofd term voor term). Het inzicht om te onthouden: de
regelmaat van het signaal en het verval van het spectrum zijn
dezelfde informatie --- een ingenieur leest de ene af van de
helling van de andere, zonder de functie ooit te tekenen.
\end{example}

\begin{theorem}[Dirichlet]\label{thm:b2:fourier:dirichlet}
Zij $f$ $2\pi$-periodiek en stuksgewijs $C^1$. Dan geldt voor elke
$x$
\[
S_N(f)(x) \xrightarrow[N \to \infty]{}
\frac{f(x^+) + f(x^-)}{2}
\]
(het gemiddelde van de eenzijdige limieten) --- in het bijzonder
$S_N(f)(x) \to f(x)$ in elk continuïteitspunt.\index{stelling van
Dirichlet}
\end{theorem}


\begin{proof}
Volgens het kernlemma en zijn gemiddelde $1$, waarbij de integraal
in de helften $u > 0$ en $u < 0$ wordt gesplitst (elk met
gemiddelde $\frac12$):
\begin{align*}
S_N(f)(x) - \frac{f(x^+) + f(x^-)}{2}
&= \frac{1}{2\pi}\int_{0}^{\pi} \bigl(f(x+u) -
f(x^+)\bigr)D_N(u)\,\dd u\\
&\quad+ \frac{1}{2\pi}\int_{-\pi}^{0}\bigl(f(x+u) -
f(x^-)\bigr)D_N(u)\,\dd u .
\end{align*}
Behandel de eerste (de tweede is \omterm{def:b2:quadratic:adjoint}{symmetrisch}). Schrijf
\[
\bigl(f(x + u) - f(x^+)\bigr)\,D_N(u)
= \underbrace{\frac{f(x+u) - f(x^+)}{\sin\frac u2}}_{g(u)}\,
\sin\Bigl(\Bigl(N + \frac12\Bigr)u\Bigr) .
\]
De functie $g$ is stuksgewijs \omterm{def:b2:metric:continuity}{continu} op $\intoc{0}{\pi}$ en heeft
een \emph{eindige limiet in $0^+$}: schrijven we
\[
g(u) = \frac{f(x+u) - f(x^+)}{u}\cdot\frac{u}{\sin\frac u2} ,
\]
dan streeft de eerste factor naar $f'(x^+)$ (eenzijdige
differentieerbaarheid, uit stuksgewijs $C^1$) en de tweede naar $2$
(de standaardlimiet $\frac{\sin v}{v} \to 1$ met $v = \frac u2$):
$g(0^+) = 2f'(x^+)$ bestaat. Dus laat $g$ zich stuksgewijs \omterm{def:b2:metric:continuity}{continu}
uitbreiden tot $\intcc{0}{\pi}$, en Riemann--Lebesgue
(\cref{thm:b2:fourier:riemannlebesgue}) stuurt de integraal naar
$0$. Dat is precies waar de hypothese voor dient: zonder
eenzijdige afgeleiden explodeert de factor $\frac{1}{\sin(u/2)}$ in
$0$ sneller dan Riemann--Lebesgue kan compenseren, en kan de
\omterm{def:b2:funcseq:def}{puntsgewijze convergentie} werkelijk falen voor louter \omterm{def:b2:metric:continuity}{continue} $f$
--- de kloof die de stelling van Fejér (weekendopgave) door
middelen dicht.
\end{proof}

\begin{example}[Dirichlet in een sprong]\label{ex:b2:fourier:jumpvalue}
Voor de zaagtand $f(t) = t$ op $\intoo{-\pi}{\pi}$
(\cref{ex:b2:fourier:basel} hieronder) springt de periodieke
uitbreiding in $t = \pi$ van $f(\pi^-) = \pi$ naar $f(\pi^+) =
-\pi$. Dirichlet belooft daar de waarde $\frac{\pi + (-\pi)}{2} =
0$, en inderdaad verdwijnt elke term van $\sum
\frac{2(-1)^{n+1}}{n}\sin nt$ in $t = \pi$: de reeks convergeert
beleefd naar het middelpunt en negeert beide eenzijdige waarden.
Verplaatsen we het evaluatiepunt naar $t = \frac\pi2$ (een
continuïteitspunt), dan maakt dezelfde reeks er Leibniz'
$\frac\pi4$ van. Eén reeks, twee gedragingen --- precies de twee
uitspraken van de stelling.
\end{example}

\begin{theorem}[Normale convergentie voor $C^1$; Parseval]\label{thm:b2:fourier:parseval}
\begin{enumerate}
  \item Is $f$ \omterm{def:b2:metric:continuity}{continu}, $2\pi$-periodiek en stuksgewijs $C^1$, dan
        is $c_n(f') = \iu n\,c_n(f)$, convergeert de \omterm{def:b2:fourier:coefficients}{Fourierreeks}
        van $f$ \emph{normaal} op $\R$, en is haar som $f$.
  \item (Parseval) Voor elke stuksgewijs \omterm{def:b2:metric:continuity}{continue}
        $2\pi$-periodieke $f$:\index{identiteit van Parseval}
        \[
        \frac{1}{2\pi}\int_{-\pi}^{\pi}\abs{f}^2
        = \sum_{n=-\infty}^{\infty} \abs{c_n(f)}^2
        = \frac{\abs{a_0}^2}{4} + \frac12\sum_{n\geq1}
        \bigl(\abs{a_n}^2 + \abs{b_n}^2\bigr).
        \]
        \emph{(Hier bewezen voor \omterm{def:b2:metric:continuity}{continue}, stuksgewijs
        $C^1$-functies $f$; in het algemeen aangenomen.)}
\end{enumerate}
\end{theorem}

\begin{proof}
(1) Partiële integratie op elk $C^1$-stuk (de randtermen heffen
elkaar op wegens \omterm{def:b2:metric:continuity}{continuïteit} en periodiciteit): $c_n(f') = \iu n
c_n(f)$. Vervolgens, met Cauchy--Schwarz op de twee
kwadratisch \omterm{def:b2:series:summable}{sommeerbare families}
(\cref{prop:b2:fourier:bessel} voor $f'$):
\[
\sum_{n \neq 0} \abs{c_n(f)} = \sum_{n\neq0}
\frac{\abs{c_n(f')}}{\abs n}
\leq \Bigl(\sum \abs{c_n(f')}^2\Bigr)^{1/2}
\Bigl(\sum_{n\neq0}\frac{1}{n^2}\Bigr)^{1/2} < \infty :
\]
\omterm{def:b2:funcseq:series}{normale convergentie} van de \omterm{def:b2:fourier:coefficients}{Fourierreeks}. Haar som is \omterm{def:b2:metric:continuity}{continu} en
valt volgens Dirichlet in elk punt met $f$ samen
(\cref{thm:b2:fourier:dirichlet}: $f$ is \omterm{def:b2:metric:continuity}{continu}): de reeks
convergeert uniform naar $f$.

(2) Voor zulke $f$: $S_N f \to f$ uniform, dus $\norm{f - S_Nf}_2
\leq \norm{f - S_Nf}_\infty \to 0$, en Pythagoras ($\norm f_2^2 =
\sum_{\abs n \leq N}\abs{c_n}^2 + \norm{f - S_Nf}_2^2$) gaat over
in de limiet. De reële vorm is boekhouding met $a_n, b_n$.
\end{proof}

\begin{example}[De $C^1$-staartschatting, kwantitatief gemaakt]\label{ex:b2:fourier:tailbound}
Het bewijs van \cref{thm:b2:fourier:parseval}~(1) verbergt een
bruikbare schatting. Voor \omterm{def:b2:metric:continuity}{continue}, stuksgewijs $C^1$-functies $f$
geeft dezelfde ongelijkheid van Cauchy--Schwarz, alleen op de
staart toegepast,
\[
\sum_{\abs n > N}\abs{c_n(f)}
= \sum_{\abs n > N}\frac{\abs{c_n(f')}}{\abs n}
\leq \Bigl(\sum_{\abs n > N}\abs{c_n(f')}^2\Bigr)^{\!1/2}
\Bigl(\sum_{\abs n>N}\frac{1}{n^2}\Bigr)^{\!1/2}
\leq \norm{f'}_2\,\sqrt{\frac{2}{N}} ,
\]
met Bessel voor $f'$ en $\sum_{n>N}n^{-2} \leq \frac1N$. De
\omterm{def:b2:funcseq:def}{uniforme} fout van de partiaalsommen voldoet dus aan
\[
\norm{f - S_Nf}_\infty
\leq \sum_{\abs n>N}\abs{c_n(f)}
\leq \norm{f'}_2\,\sqrt{\frac2N} .
\]
Voor $f(t) = \abs t$ is $\norm{f'}_2 = 1$ (de afgeleide is
$\pm1$), zodat tien termen $\abs t$ al uniform tot op $\sqrt{0.2}
\approx 0.45$ reconstrueren, en $N = 10^4$ tot op $0.015$. Het
inzicht om te onthouden: één afgeleide koopt het \omterm{def:b2:funcseq:def}{uniforme} tempo
$\frac{1}{\sqrt N}$; vergeleken met de wereld van coëfficiënten
$\frac1n$ van de blokgolf (helemaal geen \omterm{def:b2:funcseq:def}{uniforme convergentie})
krijgt het woordenboek van \cref{ex:b2:fourier:decaytable} zo
getallen.
\end{example}

\begin{example}[Basel en vrienden]\label{ex:b2:fourier:basel}
Zij $f(t) = t$ op $\intoo{-\pi}{\pi}$, $2\pi$-periodiek uitgebreid
(een zaagtand, stuksgewijs $C^1$). Rekenwerk geeft $a_n = 0$
(oneven zijn) en
\[
b_n = \frac{1}{\pi}\int_{-\pi}^{\pi} t\sin nt\,\dd t
= \frac{2(-1)^{n+1}}{n} .
\]
Dirichlet in $t = \frac\pi2$ levert Leibniz' $\frac\pi4 = 1 -
\frac13 + \frac15 - \dots$ op; Parseval geeft
\[
\frac{1}{2\pi}\int_{-\pi}^{\pi} t^2\,\dd t = \frac{\pi^2}{3}
= \frac12\sum_{n\geq1}\frac{4}{n^2}
\quad\Longrightarrow\quad
\boxed{\;\sum_{n\geq1}\frac{1}{n^2} = \frac{\pi^2}{6}\;}
\]
--- Eulers som van Basel, in twee regels. De functie $f(t) = t^2$
levert op dezelfde manier $\sum \frac1{n^4} = \frac{\pi^4}{90}$ op
(\cref{exo:b2:fourier:3}).
\end{example}

\begin{example}[Een volledige ontwikkeling met ingebouwde controle: $\abs{\sin t}$]\label{ex:b2:fourier:abssin}
De functie $f(t) = \abs{\sin t}$ is \omterm{def:b2:metric:continuity}{continu}, even, $\pi$-periodiek
(en dus $2\pi$-periodiek) en stuksgewijs $C^1$. Het even zijn doodt
de $b_n$; $a_0 = \frac1\pi\int_{-\pi}^{\pi} \abs{\sin t}\dd t =
\frac4\pi$; en voor $n \geq 1$ geeft de omzetting van een product
in een som
\[
a_n = \frac2\pi\int_0^\pi \sin t\cos nt\,\dd t
= \frac{1}{\pi}\int_0^\pi\bigl(\sin(1+n)t +
\sin(1-n)t\bigr)\dd t
= \frac2\pi\cdot\frac{1 + \cos n\pi}{1 - n^2}
\]
voor $n \neq 1$ (en $a_1 = 0$ rechtstreeks): nul voor oneven $n$,
en $a_{2k} = \frac{-4}{\pi(4k^2-1)}$. Volgens
\cref{thm:b2:fourier:parseval}~(1) is de convergentie normaal, en
geldt
\[
\abs{\sin t} = \frac{2}{\pi} - \frac{4}{\pi}
\sum_{k\geq1}\frac{\cos(2kt)}{4k^2 - 1}
\qquad (t \in \R) .
\]
Ingebouwde controle in $t = 0$: de identiteit eist dat
$\sum_{k\geq1}\frac{1}{4k^2-1} = \frac12$, wat een telescopische
som bevestigt:
\[
\sum_{k\geq1}\frac{1}{4k^2-1}
= \frac12\sum_{k\geq1}\Bigl(\frac{1}{2k-1} -
\frac{1}{2k+1}\Bigr) = \frac12 . \checkmark
\]
Het inzicht om te onthouden: het spectrum van $\abs{\sin}$ leeft
uitsluitend op de \emph{even} frequenties --- een sinus
gelijkrichten verdubbelt zijn frequentie-inhoud, en daarom brommen
dubbelfasige gelijkrichters op $100$ of $120$ hertz, tweemaal de
netfrequentie.
\end{example}

\begin{example}[Parseval als rekenmachine]\label{ex:b2:fourier:zetafourodd}
Parseval maakt van ontwikkelingen bij de vleet numerieke reeksen.
Pas haar toe op $f(t) = \abs t$ (\cref{exo:b2:fourier:2}: $a_0 =
\pi$, $a_n = \frac{-4}{\pi n^2}$ voor oneven $n$, de rest nul):
\[
\frac{1}{2\pi}\int_{-\pi}^{\pi}t^2\,\dd t = \frac{\pi^2}{3}
= \frac{a_0^2}{4} + \frac12\sum_{n \text{ oneven}} a_n^2
= \frac{\pi^2}{4} +
\frac{8}{\pi^2}\sum_{n\text{ oneven}}\frac{1}{n^4} ,
\]
waaruit
\[
\sum_{n\text{ oneven}}\frac{1}{n^4}
= \frac{\pi^2}{8}\Bigl(\frac{\pi^2}{3} -
\frac{\pi^2}{4}\Bigr) = \frac{\pi^4}{96} .
\]
Kruiscontrole met \cref{exo:b2:fourier:3}: $\sum\frac1{n^4}$ in
oneven en even delen splitsen geeft de boekhouding $S =
S_{\mathrm{oneven}} + \frac{S}{16}$ in $\zeta$-stijl, dus $S =
\frac{16}{15}\cdot\frac{\pi^4}{96} = \frac{\pi^4}{90}$ --- precies
de waarde die daar met een andere functie werd gevonden. Twee
ontwikkelingen, één getal: die overeenstemming is de isometrie van
Parseval aan het werk. Het inzicht om te onthouden: elke nieuwe
Fourierontwikkeling is een machine die reeksidentiteiten
voortbrengt; Deel II van de weekendopgave legt uit waarom de
machine zichzelf nooit kan tegenspreken.
\end{example}

\begin{omfigure}
\begin{tikzpicture}
\begin{axis}[omaxis, width=9.2cm, height=5.6cm,
    xmin=-3.6, xmax=3.6, ymin=-1.5, ymax=1.5,
    xtick={-3.1416, 0, 3.1416},
    xticklabels={$-\pi$, $0$, $\pi$}, ytick={-1,1}]
  \addplot[omExo, thick] coordinates {(-3.1416,-1) (0,-1)};
  \addplot[omExo, thick] coordinates {(0,1) (3.1416,1)};
  \addplot[omDef, thick, domain=-3.3:3.3, samples=200]
    {4/pi*sin(deg(x))};
  \addplot[omThm, very thick, domain=-3.3:3.3, samples=400]
    {4/pi*(sin(deg(x)) + sin(deg(3*x))/3 + sin(deg(5*x))/5 +
     sin(deg(7*x))/7 + sin(deg(9*x))/9)};
\end{axis}
\end{tikzpicture}

{\small De blokgolf (grijs) en de \omterm{def:b2:fourier:coefficients}{Fourierpartiaalsommen} $S_1$
(blauw) en $S_9$ (rood): convergentie in elk continuïteitspunt,
maar een blijvende overshoot van $\sim 9\%$ nabij de sprongen ---
het \emph{Gibbs-fenomeen}\index{Gibbs-fenomeen}. \omterm{def:b2:funcseq:def}{Uniforme
convergentie} faalt precies omdat de limiet discontinu is.}
\end{omfigure}

\begin{example}[Translatie en modulatie]\label{ex:b2:fourier:shiftrules}
Twee regels van één regel lang brengen uit één ontwikkeling vele
andere voort. Voor $a \in \R$ geeft de substitutie $s = t - a$
\[
c_n\bigl(f(\cdot - a)\bigr)
= \frac{1}{2\pi}\int_{-\pi}^{\pi}f(t - a)\eu^{-\iu nt}\dd t
= \eu^{-\iu na}\,c_n(f)
\qquad\text{(translatie moduleert het spectrum)},
\]
en rechtstreeks uit de definitie
\[
c_n\bigl(\eu^{\iu kt}f\bigr) = c_{n-k}(f)
\qquad\text{(modulatie transleert het spectrum)}.
\]
Uitgewerkt geval: de over $\pi$ verschoven zaagtand $g(t) = f(t -
\pi)$ met $f(t) = t$ heeft $b_n$-coëfficiënten
$(-1)^n\cdot\frac{2(-1)^{n+1}}{n} = -\frac2n$: de ontwikkeling $g
\sim -2\sum\frac{\sin nt}{n}$ van de zaagtand die in $0$ springt
in plaats van in $\pi$ --- zonder één integraal opnieuw te
berekenen. Het inzicht om te onthouden: tijdsverschuivingen draaien
alleen aan fasen, nooit aan amplituden ($\abs{c_n}$ is invariant
onder verschuiving), en daarom zijn energie (Parseval) en
convergentieklasse eigenschappen van de \emph{vorm} van het
signaal, niet van waar de klok begint.
\end{example}

\begin{remark}[Klassieke valkuilen]
\emph{(i) Drie convergenties, drie munteenheden:} puntsgewijs
(Dirichlet: vereist stuksgewijs $C^1$, betaalt in elke sprong het
\emph{middelpunt} --- nooit een eenzijdige waarde), uniform
(vereist een \omterm{def:b2:metric:continuity}{continue} limiet; onmogelijk over een sprong heen, en
Gibbs is het zichtbare symptoom) en in kwadratisch gemiddelde
(Parseval: het robuustst, blind voor afzonderlijke punten). Noem
altijd welke je beweert. \emph{(ii) Geen termsgewijze
differentiatie als vanzelf:} de zaagtandreeks van
\cref{ex:b2:fourier:basel} term voor term differentiëren levert
$\sum 2(-1)^{n+1}\cos nt$ op, waarvan de termen niet eens naar $0$
streven --- de overdrachtsstellingen van het hoofdstuk over rijen
en reeksen van functies eisen \omterm{def:b2:funcseq:def}{uniforme convergentie} van de
\emph{afgeleide} reeks, en die vernietigt de sprong. Eerst
gladheid, dan differentiëren (\cref{exo:b2:fourier:6} is het
woordenboek). \emph{(iii) \omterm{def:b2:quadratic:adjoint}{Symmetrische} partiaalsommen:} de
stelling van Dirichlet betreft $S_N = \sum_{-N}^{N}$; herschikken
of eerst één kant sommeren kan divergentie in convergentie
veranderen en terug. \emph{(iv) Verschuivende normalisatie:} de
afspraken verschillen van boek tot boek ($\frac{1}{2\pi}$ of
$\frac1\pi$ ervoor, periode $2\pi$ of $1$); de betrouwbare
invarianten zijn de orthonormaliteitsrelaties --- herbereken
$\langle e_m, e_n\rangle$ in de afspraak die voorligt voordat je
een formule vertrouwt.
\end{remark}

\begin{remark}[Waar dit wordt gebruikt]
Parseval is de kiem van de $L^2$-theorie van de \omterm{def:b2:fourier:coefficients}{Fourierreeksen}:
het volume van bachelorjaar 3 maakt het beeld af (de exponentiëlen
vormen een hilbertbasis van $L^2$, en de afbeelding $f \mapsto
(c_n)$ is een bijectieve isometrie). Binnen dit volume bewijst de
weekendopgave de \emph{stelling van Fejér} --- de
Cesàro-gemiddelden van de \omterm{def:b2:fourier:coefficients}{Fourierreeks} convergeren uniform voor
elke \omterm{def:b2:metric:continuity}{continue} periodieke $f$ --- die de aangenomen algemene
Parseval tot stelling verheft, de trigonometrische stelling van
Weierstrass oplevert en twee spectaculaire dividenden uitkeert: de
gelijkverdelingsstelling van Weyl en de isoperimetrische
ongelijkheid. De toegepaste wiskunde leest dit hoofdstuk dagelijks:
spectra van signalen, harmonischen van trillende systemen en de
snelle Fouriertransformatie (de eindige gedaante was
\cref{exo:b2:hermitian:10}).
\end{remark}

\begin{remark}[Vooruitblik binnen dit volume]
Drie hoofdstukken voeren een gesprek met dit hoofdstuk.
Terugkijkend: het \omterm{def:b2:hermitian:adjoint}{hermitische} hoofdstuk leverde de meetkunde
(orthonormale families, projecties, Bessel), en het hoofdstuk over
rijen en reeksen van functies de analyse (\omterm{def:b2:funcseq:def}{uniforme convergentie},
overdrachtsstellingen, benaderende eenheden --- de kern van Fejér
verhoudt zich tot de \omterm{def:b2:fourier:coefficients}{Fourierreeksen} als de veeltermen van
Bernstein zich tot Weierstrass verhielden). Zijwaarts: de
randtheorie van het hoofdstuk over machtreeksen keert terug via
Abelsommatie, hier verwezenlijkt door de Poisson-kern
(\cref{exo:b2:fourier:12}) --- waarbij de straal $r$ van de schijf
de rol van de sommatieparameter speelt. Vooruit: het hoofdstuk over
differentiaalvergelijkingen ontbindt een periodieke aandrijving in
harmonischen en voert elk daarvan aan de frequentierespons van de
oscillator; resonantie treedt op wanneer een Fouriermode van de
ingang met een eigenfrequentie samenvalt, en daarom zijn zijn
weekendopgave en die van dit hoofdstuk twee helften van één
verhaal.
\end{remark}

\begin{omfigure}
\begin{tikzpicture}
\begin{axis}[omaxis, width=9.2cm, height=5.6cm,
    xmin=-3.6, xmax=3.6, ymin=-2.5, ymax=9.5,
    xtick={-3.1416, 0, 3.1416},
    xticklabels={$-\pi$, $0$, $\pi$}, ytick={8}]
  \addplot[omDef, thick, domain=0.02:3.3, samples=300]
    {sin(deg(8.5*x))/sin(deg(0.5*x))};
  \addplot[omDef, thick, domain=-3.3:-0.02, samples=300]
    {sin(deg(8.5*x))/sin(deg(0.5*x))};
  \addplot[omThm, very thick, domain=0.02:3.3, samples=300]
    {sin(deg(4*x))^2/(8*sin(deg(0.5*x))^2)};
  \addplot[omThm, very thick, domain=-3.3:-0.02, samples=300]
    {sin(deg(4*x))^2/(8*sin(deg(0.5*x))^2)};
  \node[font=\small, omDef] at (axis cs:1.05,5.2) {$D_8$};
  \node[font=\small, omThm] at (axis cs:-0.75,6.2) {$F_8$};
\end{axis}
\end{tikzpicture}

{\small De \omterm{lem:b2:fourier:kernel}{Dirichlet-kern} $D_8$ (blauw) oscilleert en neemt
negatieve waarden aan; de kern $F_8$ van Fejér (rood) is
niet-negatief, concentreert zich in $0$ en heeft gemiddelde $1$:
een \emph{benaderende eenheid}. Positiviteit is precies wat de
\omterm{lem:b2:fourier:kernel}{Dirichlet-kern} mist, en wat de stelling van Fejér in de
weekendopgave onvoorwaardelijk maakt.}
\end{omfigure}

\section{Oefeningen}

\begin{exercise}[$\star$]\label{exo:b2:fourier:1}
Bereken de \omterm{def:b2:fourier:coefficients}{Fouriercoëfficiënten} van de blokgolf ($f = -1$ op
$\intoo{-\pi}{0}$, $+1$ op $\intoo{0}{\pi}$), formuleer het besluit
van Dirichlet in $t = \frac\pi2$ en in de sprong $t = 0$, en vind
de reeks van Leibniz terug.
\end{exercise}

\begin{exercise}[$\star$]\label{exo:b2:fourier:2}
Ontwikkel $f(t) = \abs t$ ($\abs t \leq \pi$, $2\pi$-periodiek) in
\omterm{def:b2:fourier:coefficients}{Fourierreeks}; verantwoord de \emph{normale} convergentie; evalueer
in $t = 0$ om $\sum_{k\geq0}\frac{1}{(2k+1)^2} = \frac{\pi^2}{8}$
te verkrijgen, en leid daaruit Basel opnieuw af.
\end{exercise}

\begin{exercise}[$\star$]\label{exo:b2:fourier:3}
Ontwikkel $f(t) = t^2$ ($\abs t \leq \pi$) en leid af dat
\[
\sum_{n\geq1}\frac{(-1)^{n+1}}{n^2} = \frac{\pi^2}{12},
\qquad
\sum_{n\geq1}\frac{1}{n^4} = \frac{\pi^4}{90}
\quad\text{(Parseval)}.
\]
\end{exercise}

\begin{exercise}[$\star\star$]\label{exo:b2:fourier:4}
Zij $f$ \omterm{def:b2:metric:continuity}{continu} en $2\pi$-periodiek met $c_n(f) = 0$ voor alle $n$.
Bewijs dat $f = 0$ \emph{(Parseval --- voor welke klasse is zij
hier bewezen? verantwoord dat \omterm{def:b2:metric:continuity}{continuïteit} plus stuksgewijs $C^1$
mag vervallen door de algemene Parseval aan te nemen, of geef het
dichtheidsargument in grote lijnen)}.
\end{exercise}

\begin{exercise}[$\star\star$]\label{exo:b2:fourier:5}
Ontwikkel voor $\alpha \notin \Z$ de functie $f(t) = \cos(\alpha
t)$ ($\abs t \leq \pi$) en leid de splitsing in partieelbreuken van
de cotangens af:
\[
\pi\cot(\pi\alpha) = \frac{1}{\alpha} + \sum_{n\geq1}
\frac{2\alpha}{\alpha^2 - n^2} .
\]
\end{exercise}

\begin{exercise}[$\star\star$]\label{exo:b2:fourier:6}
Bewijs dat als $f$ $2\pi$-periodiek en van klasse $C^k$ is met
$f^{(k)}$ stuksgewijs \omterm{def:b2:metric:continuity}{continu}, dan $c_n(f) = O\bigl(\abs
n^{-k}\bigr)$: gladheid van het signaal $=$ verval van zijn
spectrum.
\end{exercise}

\begin{exercise}[$\star\star\star$]\label{exo:b2:fourier:7}
(Ongelijkheid van Wirtinger) Zij $f$ van klasse $C^1$,
$2\pi$-periodiek, met $\int_{-\pi}^{\pi} f = 0$. Bewijs dat
\[
\int_{-\pi}^{\pi} \abs{f}^2 \leq \int_{-\pi}^{\pi} \abs{f'}^2 ,
\]
met gelijkheid dan en slechts dan als $f(t) = a\cos t + b\sin t$.
\emph{(Parseval op beide leden; vergelijk $\abs{c_n}^2$ en
$n^2\abs{c_n}^2$.)}
\end{exercise}

\begin{exercise}[$\star\star\star$]\label{exo:b2:fourier:8}
(De constante van Gibbs) Evalueer voor de blokgolf van
\cref{exo:b2:fourier:1} de partiaalsom in $x_N = \frac{\pi}{2N}$:
toon aan, met $u_k = \frac{(2k+1)\pi}{2N}$ en $\Delta u =
\frac{\pi}{N}$, dat
\[
S_{2N-1}\Bigl(\frac{\pi}{2N}\Bigr)
= \frac{4}{\pi}\sum_{k=0}^{N-1} \frac{\sin u_k}{2k+1}
= \frac{2}{\pi}\sum_{k=0}^{N-1} \frac{\sin u_k}{u_k}\,\Delta u
\xrightarrow[N\to\infty]{}
\frac{2}{\pi}\int_0^{\pi}\frac{\sin u}{u}\,\dd u \approx 1.179 :
\]
een Riemannsom van $\frac{2}{\pi}\cdot\frac{\sin u}{u}$ op
$\intcc{0}{\pi}$ in de middelpunten. Besluit dat de overshoot
boven de sprongwaarde $1$ niet verdwijnt als $N \to \infty$.
\end{exercise}

\begin{exercise}[$\star$]\label{exo:b2:fourier:9}
Ontwikkel $\cos^3 t$ en $\sin^2 t\,\cos t$ in \omterm{def:b2:fourier:coefficients}{Fourierreeks}
\emph{(lineariseer; een trigonometrische veelterm is haar eigen
\omterm{def:b2:fourier:coefficients}{Fourierreeks}, wegens de eenduidigheid van de coëfficiënten)}. Wat
zijn $c_n$, $a_n$ en $b_n$ in elk geval?
\end{exercise}

\begin{exercise}[$\star\star$]\label{exo:b2:fourier:10}
Zij $a > 0$ en $f(t) = \eu^{at}$ op $\intoc{-\pi}{\pi}$,
$2\pi$-periodiek uitgebreid. Bereken
\[
c_n(f) = \frac{(-1)^n\sinh(a\pi)}{\pi(a - \iu n)} ,
\]
pas de stelling van Dirichlet toe in de sprong $t = \pi$, en leid
de splitsing in partieelbreuken van de hyperbolische cotangens af:
\[
\coth(\pi a) = \frac{1}{\pi a} + \sum_{n\geq1}
\frac{2a}{\pi(a^2 + n^2)} .
\]
\end{exercise}

\begin{exercise}[$\star\star$]\label{exo:b2:fourier:11}
(Convolutie) Definieer voor \omterm{def:b2:metric:continuity}{continue}, $2\pi$-periodieke $f, g$
\[
(f * g)(x) = \frac{1}{2\pi}\int_{-\pi}^{\pi} f(x - t)\,g(t)\,
\dd t .
\]
Toon aan dat $f * g = g * f$, dat $c_n(f*g) = c_n(f)\,c_n(g)$
\emph{(om de twee integralen van een \omterm{def:b2:metric:continuity}{continue} integrand te
verwisselen, vergelijk de twee functies van de bovengrens: beide
verdwijnen in het linkerrandpunt en hebben dezelfde afgeleide,
wegens \omterm{def:b2:metric:continuity}{continuïteit} en differentiatie onder het integraalteken)},
en dat $S_N(f) = f * D_N$ voor de \omterm{lem:b2:fourier:kernel}{Dirichlet-kern}. (De
Fejér-gemiddelden van de weekendopgave zijn eveneens convoluties,
$\sigma_N(f) = f * F_N$.)
\end{exercise}

\begin{exercise}[$\star\star\star$]\label{exo:b2:fourier:12}
(Poisson-kern: Abelgemiddelden van \omterm{def:b2:fourier:coefficients}{Fourierreeksen}) Zet voor $0 \leq
r < 1$ de functie $P_r(t) = \sum_{n\in\Z} r^{\abs n}\eu^{\iu nt}$.
\begin{enumerate}
  \item Sommeer de twee meetkundige reeksen en toon aan dat
        \[
        P_r(t) = \frac{1 - r^2}{1 - 2r\cos t + r^2} > 0,
        \qquad
        \frac{1}{2\pi}\int_{-\pi}^{\pi}P_r = 1 .
        \]
  \item Toon aan dat voor $\delta \leq \abs t \leq \pi$ geldt
        $P_r(t) \leq \frac{1 - r^2}{1 - 2r\cos\delta + r^2} \to 0$
        als $r \to 1^-$, uniform.
  \item Leid af dat voor elke \omterm{def:b2:metric:continuity}{continue}, $2\pi$-periodieke $f$ de
        \emph{Abelgemiddelden} $(f * P_r)(x) = \sum_n r^{\abs
        n}c_n(f)\,\eu^{\iu nx}$ uniform naar $f$ convergeren als
        $r \to 1^-$ --- de zusterversie met \omterm{def:b2:metric:continuity}{continue} parameter van
        de stelling van Fejér, en de Fouriergedaante van de
        Abelsommatie uit het hoofdstuk over machtreeksen.
\end{enumerate}
\end{exercise}

\section{Probleem: de stelling van Fejér en haar dividenden}

\begin{problem}\label{pb:b2:fourier:1}
De stelling van Dirichlet heeft $f$ stuksgewijs $C^1$ nodig; voor
louter \omterm{def:b2:metric:continuity}{continue} $f$ kunnen de partiaalsommen $S_N(f)$ zich
misdragen. Fejérs ontdekking: hun \emph{Cesàro-gemiddelden} nooit.
De motor is de positiviteit van de kern van Fejér, en de oogst is
enorm: \omterm{def:b2:funcseq:def}{uniforme} trigonometrische benadering (Weierstrass),
eenduidigheid van de \omterm{def:b2:fourier:coefficients}{Fouriercoëfficiënten}, Parseval voor elke
functie van dit hoofdstuk (waarmee het ``aangenomen'' in
\cref{thm:b2:fourier:parseval} verdwijnt), de
gelijkverdelingsstelling van Weyl, en --- als kroon op een eeuw
meetkunde --- de isoperimetrische ongelijkheid.\index{stelling van
Fejer@stelling van Fejér}\index{gelijkverdeling}\index{isoperimetrische
ongelijkheid} Overal is $f$ $2\pi$-periodiek en stuksgewijs
\omterm{def:b2:metric:continuity}{continu}, en
\[
\sigma_N(f) = \frac{S_0(f) + S_1(f) + \dots +
S_{N-1}(f)}{N} .
\]

\medskip
\noindent\textbf{Deel I --- De kern van Fejér.}
\begin{enumerate}
  \item Toon aan dat $\sigma_N(f)(x) =
        \frac{1}{2\pi}\int_{-\pi}^{\pi}f(x+u)\,F_N(u)\,\dd u$ met
        $F_N = \frac{D_0 + \dots + D_{N-1}}{N}$, en dat
        $\frac{1}{2\pi}\int_{-\pi}^{\pi}F_N = 1$.
  \item Bewijs de gesloten vorm, voor $u \notin 2\pi\Z$:
        \[
        F_N(u) = \frac{1}{N}\,
        \frac{\sin^2\bigl(\frac{Nu}{2}\bigr)}
        {\sin^2\bigl(\frac u2\bigr)} \;\geq\; 0
        \]
        \emph{(sommeer $\sin\bigl((n+\frac12)u\bigr)$ als het
        imaginaire deel van een meetkundige reeks)}.
  \item Toon de concentratieschatting aan: voor $0 < \delta \leq
        \abs u \leq \pi$ is
        \[
        F_N(u) \leq \frac{1}{N\sin^2\frac\delta2}
        \xrightarrow[N\to\infty]{} 0
        \quad\text{uniform} :
        \]
        $(F_N)$ is een positieve benaderende eenheid.
  \item (Stelling van Fejér) Bewijs: is $f$ \omterm{def:b2:metric:continuity}{continu} en
        $2\pi$-periodiek, dan $\sigma_N(f) \to f$
        \emph{uniform} op $\R$ \emph{(splits de integraal van
        $\bigl(f(x+u) - f(x)\bigr)F_N(u)$ bij $\abs u = \delta$;
        gebruik Heine en de vragen 1--3)}.
  \item Toon voor stuksgewijs \omterm{def:b2:metric:continuity}{continue} $f$ de \omterm{def:b2:funcseq:def}{puntsgewijze} versie
        $\sigma_N(f)(x) \to \frac{f(x^+) + f(x^-)}{2}$ in elke $x$
        aan, en de \omterm{def:b2:funcseq:def}{uniforme} grens $\norm{\sigma_N(f)}_\infty \leq
        \norm f_\infty$ \emph{(positiviteit!)}.
\end{enumerate}

\medskip
\noindent\textbf{Deel II --- Weierstrass, eenduidigheid,
Parseval.}
\begin{enumerate}[resume]
  \item (Trigonometrische Weierstrass) Leid af: elke \omterm{def:b2:metric:continuity}{continue}
        $2\pi$-periodieke functie is een \omterm{def:b2:funcseq:def}{uniforme} limiet van
        trigonometrische veeltermen.
  \item (Eenduidigheid) Leid af: een \omterm{def:b2:metric:continuity}{continue} $f$ met $c_n(f) = 0$
        voor alle $n$ is identiek nul --- twee \omterm{def:b2:metric:continuity}{continue}
        periodieke functies met dezelfde \omterm{def:b2:fourier:coefficients}{Fouriercoëfficiënten}
        vallen samen (\cref{exo:b2:fourier:4}, nu zonder iets aan
        te nemen).
  \item (Parseval, \omterm{def:b2:metric:continuity}{continu} geval) Bewijs met de
        projectie-eigenschap van $S_N$
        (\cref{prop:b2:fourier:bessel}) en $\sigma_Nf \in \mathcal
        T_{N-1} \subseteq \mathcal T_N$ dat
        \[
        \norm{f - S_Nf}_2 \leq \norm{f - \sigma_Nf}_2
        \leq \norm{f - \sigma_Nf}_\infty
        \xrightarrow[N\to\infty]{} 0 ,
        \]
        en besluit tot de identiteit van Parseval voor elke
        \emph{\omterm{def:b2:metric:continuity}{continue}} $2\pi$-periodieke $f$.
  \item (Parseval, stuksgewijs \omterm{def:b2:metric:continuity}{continu} geval) Construeer, gegeven
        stuksgewijs \omterm{def:b2:metric:continuity}{continue} $f$ en $\varepsilon > 0$, een
        \omterm{def:b2:metric:continuity}{continue} periodieke $g$ met $\norm{f - g}_2 \leq
        \varepsilon$ \emph{(vervang $f$ door een affiene
        interpolatie op piepkleine intervallen rond de
        sprongen)}, en leid af dat $\norm{f - S_Nf}_2 \to 0$
        \emph{(gebruik Bessel: $\norm{S_Nh}_2 \leq \norm h_2$)}:
        Parseval geldt in de volle algemeenheid van
        \cref{thm:b2:fourier:parseval} --- het ``aangenomen'' is
        verdwenen.
  \item (Geen Gibbs bij Fejér) Zet dit af tegen
        \cref{exo:b2:fourier:8}: toon aan dat voor de blokgolf $f$
        geldt $\abs{\sigma_N(f)} \leq 1$ overal en voor elke $N$
        --- middelen volgens Cesàro wist de overshoot uit die
        $S_N$ achtervolgt. Leg in één zin uit welke eigenschap van
        $F_N$ hiervoor verantwoordelijk is.
\end{enumerate}

\medskip
\noindent\textbf{Deel III --- Snelheden.}
\begin{enumerate}[resume]
  \item Bewijs de twee grenzen voor de kern, voor $0 < \abs u \leq
        \pi$:
        \[
        F_N(u) \leq N,
        \qquad
        F_N(u) \leq \frac{\pi^2}{N u^2}
        \]
        \emph{(voor de eerste: $\abs{\sin N\theta} \leq
        N\abs{\sin\theta}$ met inductie; voor de tweede:
        $\sin\frac u2 \geq \frac{u}{\pi}$ op $\intcc{0}{\pi}$)}.
  \item Leid de schatting van het eerste moment af,
        \[
        \frac{1}{2\pi}\int_{-\pi}^{\pi}\abs u\,F_N(u)\,\dd u
        \;\leq\; \frac{C\,\ln N}{N}
        \qquad (N \geq 2)
        \]
        voor een expliciete constante \emph{(splits bij $\abs u =
        \frac1N$)}.
  \item Besluit: is $f$ $L$-lipschitz en $2\pi$-periodiek, dan is
        \[
        \norm{\sigma_N f - f}_\infty \leq \frac{C\,L\ln N}{N} .
        \]
  \item (Verzadiging) Bereken $\sigma_N(e_1)$ voor $e_1(t) =
        \eu^{\iu t}$ en toon aan dat $\norm{\sigma_Ne_1 -
        e_1}_\infty = \frac1N$: zelfs voor de gladste functies
        convergeert Fejér niet sneller dan $\frac1N$ --- het
        exacte analogon van de verzadiging van Bernstein in de
        weekendopgave van het hoofdstuk over rijen en reeksen van
        functies.
  \item (Lokalisatie) Toon aan: verdwijnt $f$ (stuksgewijs
        \omterm{def:b2:metric:continuity}{continu}) op $\intoo{x - \delta}{x + \delta}$, dan is
        $\sigma_N(f)(x) \to 0$, hoe wild $f$ elders ook is --- de
        convergentie van de gemiddelden in $x$ ziet alleen $f$
        nabij $x$.
\end{enumerate}

\medskip
\noindent\textbf{Deel IV --- De gelijkverdelingsstelling van
Weyl.} Een rij $(x_n)_{n\geq1}$ in $\intco{0}{1}$ heet
\emph{gelijkverdeeld} wanneer voor elk interval $\intcc{a}{b}
\subseteq \intcc{0}{1}$ geldt
\[
\frac{\#\{n \leq N : x_n \in \intcc ab\}}{N}
\xrightarrow[N\to\infty]{} b - a .
\]
\begin{enumerate}[resume]
  \item Toon aan dat $(x_n)$ gelijkverdeeld is zodra
        $\frac1N\sum_{n\leq N}f(x_n) \to \int_0^1 f$ voor elke
        \emph{\omterm{def:b2:metric:continuity}{continue}} $1$-periodieke $f$ \emph{(klem de
        indicator van $\intcc ab$ tussen twee \omterm{def:b2:metric:continuity}{continue},
        stuksgewijs affiene functies waarvan de integralen $\varepsilon$
        verschillen)}.
  \item (Criterium van Weyl, voldoende voorwaarde) Neem aan dat
        \[
        \frac{1}{N}\sum_{n=1}^{N}\eu^{2\iu\pi kx_n}
        \xrightarrow[N\to\infty]{} 0
        \qquad\text{voor elke } k \in \Z\setminus\{0\} .
        \]
        Toon aan dat $\frac1N\sum f(x_n) \to \int_0^1f$, eerst
        voor trigonometrische veeltermen, dan voor alle \omterm{def:b2:metric:continuity}{continue}
        $1$-periodieke $f$ dankzij vraag 6 (overgezet naar periode
        $1$): met vraag 16 is $(x_n)$ gelijkverdeeld.
  \item Zij $\alpha$ irrationaal en $x_n = \{n\alpha\}$ (het
        gebroken deel). Begrens de meetkundige som
        \[
        \Bigl|\sum_{n=1}^{N}\eu^{2\iu\pi kn\alpha}\Bigr|
        \leq \frac{2}{\abs{1 - \eu^{2\iu\pi k\alpha}}}
        \qquad (k \neq 0),
        \]
        en besluit tot de \emph{stelling van Weyl}:
        $(\{n\alpha\})$ is gelijkverdeeld in $\intco{0}{1}$.
  \item Leid af dat $(\{n\alpha\})$ dicht ligt in $\intcc{0}{1}$
        voor irrationale $\alpha$, en leg in één zin uit waarom
        gelijkverdeling strikt sterker is dan dichtheid.
  \item (Begincijfers) Bewijs dat het aandeel van de gehele
        getallen $n \leq N$ waarvoor $2^n$ als eerste (decimale)
        cijfer een $1$ heeft, naar $\log_{10}2 \approx 0.301$
        streeft \emph{(begincijfer $1$ betekent
        $\{n\log_{10}2\} \in \intco{0}{\log_{10}2}$; toon aan dat
        $\log_{10}2$ irrationaal is)}.
\end{enumerate}

\medskip
\noindent\textbf{Deel V --- De isoperimetrische ongelijkheid.} Zij
$\Gamma$ een gesloten enkelvoudige $C^1$-kromme van lengte $L$ die
een georiënteerde oppervlakte $A$ omsluit, geparametriseerd naar
geschaalde booglengte: $z(t) = x(t) + \iu y(t)$, $2\pi$-periodiek,
met $\abs{z'(t)} = \frac{L}{2\pi}$ constant; de omsloten
oppervlakte is
\[
A = \frac12\int_0^{2\pi}\bigl(x\,y' - y\,x'\bigr)\dd t
= \frac{1}{2}\,\Im\int_0^{2\pi}\conj{z}\,z'\,\dd t
\]
(hier als definitie van de georiënteerde oppervlakte genomen; het
hoofdstuk over meervoudige integralen bewijst dat zij met de
intuïtieve overeenstemt, via de formule van Green).
\begin{enumerate}[resume]
  \item Ontwikkel $z(t) = \sum_{n\in\Z}c_n\eu^{\iu nt}$ (de reeks
        van een $C^1$-functie, normaal convergent) en bewijs met
        Parseval toegepast op $z'$:
        \[
        \frac{L^2}{2\pi} = \int_0^{2\pi}\abs{z'}^2\dd t
        = 2\pi\sum_{n\in\Z}n^2\abs{c_n}^2 .
        \]
  \item Bewijs evenzo dat $A = \pi\sum_{n\in\Z} n\,\abs{c_n}^2$
        \emph{(Parseval in gepolariseerde vorm:
        $\frac{1}{2\pi}\int\conj f g = \sum \conj{c_n(f)}c_n(g)$,
        toegepast op $f = z$, $g = z'$)}.
  \item (Hurwitz) Besluit:
        \[
        L^2 - 4\pi A = 4\pi^2\sum_{n\in\Z}
        (n^2 - n)\abs{c_n}^2 \;\geq\; 0 ,
        \]
        met gelijkheid dan en slechts dan als $z(t) = c_0 +
        c_1\eu^{\iu t}$ --- een cirkel. De isoperimetrische
        ongelijkheid: van alle gesloten krommen van lengte $L$
        omsluit alleen de cirkel de oppervlakte
        $\frac{L^2}{4\pi}$.
  \item Verstandscontroles: ga de gelijkheid na voor de cirkel met
        straal $R$ en de strikte ongelijkheid voor het vierkant
        met zijde $a$; leg uit waarom $n^2 - n \geq 0$ voor elk
        geheel getal $n$, ook de negatieve, en waar de constante
        snelheid van de parametrisering werd gebruikt.
  \item Synthese. In telkens één zin: (i) de ene eigenschap van
        $F_N$ waaruit de Delen I--III voortvloeien en die $D_N$
        mist; (ii) hoe de sommatie volgens Cesàro hier samenhangt
        met de weekendopgave van het hoofdstuk over machtreeksen
        (Frobenius); (iii) welk dividend alleen Weierstrass (vraag
        6) gebruikte en welk de volledige Parseval nodig had; (iv)
        één zin over wat het volume van bachelorjaar 3 toevoegt
        (de \omterm{def:b2:metric:complete}{volledigheid} van $L^2$: \omterm{def:b2:fourier:coefficients}{Fourierreeksen} als
        hilbertbasis).
\end{enumerate}
\end{problem}
